\subsection{Stronger lower bounds from formal groups}

\begin{proposition}[B-RD18: every admissible ramification has a height-two counterexample]
\label{prop:B-RD18}
\status{Proved; independent construction and determinant audit passed.}
Let $p$ be any prime, let $e$ be a positive multiple of $p-1$, and
put $s=v_p(e/(p-1))$. Suppose $n\geq1$ satisfies
\[
 s\geq n-1,\qquad
 e\geq(p-1)p^{\max(2n-2,\,4n-4-s)}.                \tag{B-RD18.1}
\]
Then there are a complete discretely valued field $K$ with
algebraically closed residue field of characteristic $p$, absolute
ramification $e$, and $\mu_p\subset K$, and two height-two BT groups
$M,N$ over $K$ with
\[
 M[p^n]\simeq N[p^n],
\]
whose characters on prime-to-$p$ inertia differ. Both groups
individually acquire good reduction. The group $M$ can be taken to
have tame good-reduction defect $p+1$, and $N$ to be constant etale.
The same examples in every height $h\geq2$ are obtained by adding
constant etale groups.
\end{proposition}
\begin{proof}
Choose a height-two Honda formal group over
$k=\overline{\mathbf F}_p$. Its endomorphism ring contains
$W(\mathbf F_{p^2})$, so the Teichmuller units give a tame
order-$p+1$ automorphism $\alpha$ whose tangent eigenvalue $\zeta$
is primitive of that order. The two eigenvalues on its isocrystal
are $\zeta,\zeta^p=\zeta^{-1}$.

Apply the full-level Drinfeld deformation-ring construction in the
proof of Theorem~\ref{thm:B-RD12}, with $m=p+1$. Every ingredient
before the elliptic algebraization step involves only the formal
height-two group: the regular full-level ring, the two torsion
coordinates, averaging for the tame action, the product formula
for $p$, its valuation regions, and the equivariant discrete
valuation specialization. The condition $m\mid p+1$ holds, and its
arithmetic existence condition is exactly (B-RD18.1), because
$(p^2-1)/m=p-1$.

Stop at that formal-group construction and descend its generic BT
group along the cyclic tame extension. Over a characteristic-zero
field this is ordinary Galois descent for the compatible finite
etale groups at all levels. The full level structure is fixed by
the descent, so $M[p^n]$ is constant over $K$. The special descent
action is $\alpha$, hence is faithful of order $p+1$. Its trace at a
generator is $\zeta+\zeta^{-1}\ne2$.
The constant etale height-two group $N$ has the same constant
$p^n$-torsion and has inertia trace two. Adding constant etale
height $h-2$ adds $h-2$ to both traces and preserves their difference.

The proof also applies at $p=2$: the chosen order is three, so every
averaging denominator remains a unit; all Drinfeld product and
valuation formulas used in the construction hold for that prime.
Finally the cyclotomic containment can be checked directly, without
assuming a polarization of $M$. Over the good extension, the
rank-one crystalline representation $\det(T_pM)\otimes\mathbf Q_p$
has the single Hodge--Tate weight one. Kisin's
Corollary~2.2.6 realizes it by a height-one good group, which over
this strictly henselian complete base is multiplicative. Thus
$\det(T_pM)\chi_p^{-1}$ is a finite character. Its inertia type is
$\det(\alpha)=\zeta^{p+1}=1$, so it is unramified, hence trivial
because the residue field is algebraically closed. Therefore
$\det(T_pM)=\chi_p$. The full rational $p^n$-torsion of $M$ forces
$\chi_p\equiv1\pmod{p^n}$, proving $\mu_{p^n}\subset K$.
\end{proof}

\begin{corollary}[An explicit general-height lower bound]
\label{cor:B-RD18}
Let $h\geq2$, $\mu_p\subset K$, and define
\[
 L=\left\lfloor\log_p\!\left(\frac e{p-1}\right)\right\rfloor,
 \qquad
 N=1+\min\left\{s,\left\lfloor\frac L2\right\rfloor,
                         \left\lfloor\frac{L+s}4\right\rfloor\right\}.
\]
Any universal least index in B.4 satisfies $\nu\geq N$.
In particular $N\geq1$, so for every admissible $e$ and every
$h\geq2$, the first truncation $M(0)$ does not determine the
requested character. This is stronger than restricting the lower
bound to elliptic examples: the tame automorphism of a formal
height-two group need not arise from an automorphism of an
elliptic curve.
\end{corollary}
\begin{proof}
The integer equivalence between (B-RD18.1) and $n\leq N$ is the
same calculation as in Corollary~\ref{cor:B-RD13}, with
$m=p+1$. Failure at level $p^N=M(N-1)$ gives $\nu\geq N$.
\end{proof}

\begin{remark}[Finite arithmetic certificate]
The proof above is symbolic. As a separate finite check,
\path{code/verify_honda_bounds.py} enumerates the exact integer-pair
formula for $p=2,3,5,7,11,13,17,19$, levels $n=1,2,3$, and
$e\leq1200$, always with $m=p+1$. It verifies all 24 range tables,
3803 constructed pairs, and 11088 least-level formula comparisons
through $e\leq5000$. All checks pass; the recorded run took
$0.508$ seconds. This certificate verifies the integer arithmetic,
not the formal deformation theory.
\end{remark}
